\documentclass[11pt]{article}
\usepackage{fontspec}
\usepackage[french,english]{babel}
\usepackage{amsmath}
\usepackage{amssymb}
\usepackage{amsthm}
\usepackage{mathrsfs}
\usepackage{tikz-cd}
\usepackage{enumitem}

\newcommand{\printedpage}[1]{\marginpar{\footnotesize p.~#1}}
\newcommand{\msp}[1]{\marginpar{\footnotesize [#1]}}
\newcommand{\card}{\operatorname{card}}

\begin{document}

\printedpage{174}
% head: A. Grothendieck : Dérivateurs, Ch. XVIII, Édition des Universités de Montpellier II et Paris VII

\msp{page 451}
Saying that it is the identity, is expressed by condition 2°) on $(x_j,p_j)_{j\in J}$. Condition 3° is clear by construction, $p_j$ admitting an image isomorphic to $z$ via $\alpha_j$. On the other hand, one sees \texttt{[right]} away that for an arrow $\lambda : j \longrightarrow j'$ of $J$, one has commutativity in % [sic] est clair -> est claire, agreement with "condition"
\[
\begin{tikzcd}
x_j \arrow[r, "x_\lambda"] & x_{j'} \\
x_j^\natural \arrow[u, hook] \arrow[r, "x_\lambda"] & x_{j'}^\natural \arrow[u, hook] \\
z \arrow[u, "\wr", "\overline{\alpha_j}"'] \arrow[r, "\mathrm{id}_z", "\sim"'] & z \ , \arrow[u, "\wr", "\overline{\alpha_{j'}}"']
\end{tikzcd}
\]
where $\overline{\alpha_j}$, $\overline{\alpha_{j'}}$ are the isomorphisms defined by $\alpha_j$, $\alpha_{j'}$ respectively. This proves that $x_\lambda^\natural$ is an isomorphism.

This proves that if $(x_i)_{i\in I}$ is essentially constant, there exists $J\longrightarrow I$ as indicated. The converse is also true, for one can suppose $J = I$, since $(x_i)_{i\in I}$ and $(x_j)_{j\in J}$ define isomorphic ind-objects. Then the direct system of the $x_j^\natural$ has as transition morphisms isomorphisms, and as $J$ is filtering hence 1-connected, one easily deduces from this that it is isomorphic
\msp{page 452}
to a constant direct system, with value $z$. Choosing such an isomorphism
\[
z \xrightarrow[\sim]{\overline{\alpha_i}} x_i^\natural \ ,
\]
one finds $z \xrightarrow{\alpha} \underline{x}$, and the system of the $p_i$ defines some
\[
\begin{tikzcd}
x_i \arrow[r, "\overline{p_i}"] \arrow[rr, bend right=12, "\beta_i \overset{\text{def}}{=} \overline{\alpha_j}^{-1}\overline{p_i}"'] & x_i^\natural \arrow[r, "\overline{\alpha_j}^{-1}"] & z \ ,
\end{tikzcd}
\] % [sic] \overline{\alpha_j} (index j) -> \overline{\alpha_i}, the arrow goes from x_i^\natural, cf. z -> x_i^\natural above
and one easily sees that one will have
\[
\beta\alpha \ = \ id_z \ , \qquad \alpha\beta \ = \ \mathrm{id}_{\underline{x}} \ ,
\]

\hfill q.e.d.

\textsc{I can now prove} 4.15.3. It is indeed essentially the conjunction of 4.15.7.5 (p. 448), and of 4.8.13, which tells us essentially that if $\mathcal{C}$ is \textit{Karoubian} (a condition satisfied if $\mathcal{C}$ is without projectors), then every direct system in $\mathcal{C}$, large relative to $\pi_0$, where $\pi_0 \geqslant \card\,\text{red}\,\mathcal{C}$, is essentially constant,

\hfill q.e.d.

\printedpage{175}
% head: A. Grothendieck : Dérivateurs, Ch. XVIII, Édition des Universités de Montpellier II et Paris VII

\msp{page 453}
\textbf{4.15.8.} I want to return here to the unfinished proof of 4.8.13 (p. 212-216). I take up the proof where I had left it (p. 215). One must prove that
\[
F(\xi) \ = \ \varinjlim_{i\in I} M_i
\]
has cardinality $\leqslant \pi_0$. I announced that one can reduce to the case where $M$ is a monoid of monomorphisms, i.e. that $uv = uv' \Longrightarrow v = v'$ (for $u,v,v' \in M$). I am going to make this argument explicit.

For every $i\in I$, let
\[
N_i \ = \ \left\{ u \in M \ \middle| \ \begin{array}{l}
\text{the set } I_{i,u} \text{ of } j\in I \text{ such that}\\
j \geqslant i,\ \alpha_{ji} = u \text{ is cofinal}\\
\text{in } I\text{, i.e. for every } j_0\in I,\\
\text{there exists } j \in I_{i,u}\text{, with } j \geqslant j_0
\end{array} \right\} .
\tag*{(4.15.8.1)}
\]
Since $I$, hence $i\backslash I$ is large relative to $\pi_0$, and $\card M \leqslant \pi_0$, one sees at once that
\[
N_i \ \neq \ \emptyset \ .
\tag*{(4.15.8.2)}
\]
$\Big($ If on the other hand $i' \geqslant i$, one finds
\[
N_i \ = \ N_{i'}\alpha_{i'i} \ .
\tag*{(4.15.8.3)}
\]
The inclusion $N_{i'}\alpha_{i'i} \subset N_i$ is $\pm$ tautological, let us prove $N_i \subset N_{i'}\alpha_{i'i}$. Let $u \in N_i$, the set $J = I_{i,u} \cap (i'\backslash I)$ of $j \geqslant i'$ such that $\alpha_{ji} = u$ is cofinal in $I$, hence large relative to $\pi_0$.
\[
\begin{tikzcd}
i \arrow[r] \arrow[rr, bend right=12, "u = \alpha_{ji}"'] & i' \arrow[r, "\alpha_{ji'}"] & j
\end{tikzcd}
\]
Thus there exists a $v \in M$ such that the set $J'$ of $j\in J$ such that $\alpha_{ji'} = v$
\msp{page 454}
is cofinal in $J$, hence in $I$. Since this set $J'$ is contained in $I_{i',v}$, this implies that $I_{i',v}$ is cofinal in $I$, i.e. $v \in N_{i'}$. On the other hand one has, for $j \in J'$,
\[
\underbrace{\alpha_{ji'}}_{v}\alpha_{i'i} \ = \ \underbrace{\alpha_{ji}}_{u} \ ,
\]
whence $u = v\alpha_{i'i}$ with $v\in N_{i'}$, whence $u \in N_i\alpha_{i'i}$, q.e.d. $\Big)$ % [sic] N_i\alpha_{i'i} -> N_{i'}\alpha_{i'i}, the inclusion to be proved is N_i \subset N_{i'}\alpha_{i'i}

For $i \in I$, let
\[
\alpha_i \ : \ \underbrace{\mathrm{Hom}(\xi, \underbrace{x_i}_{=\,\xi})}_{=\,M} \ \longrightarrow \ \underbrace{\mathrm{Hom}(\xi,F) \ = \ F(\xi)}_{\simeq\ \varinjlim\limits_{i} \mathrm{Hom}(\xi,x_i)}
\]
\printedpage{176}
% head: A. Grothendieck : Dérivateurs, Ch. XVIII, Édition des Universités de Montpellier II et Paris VII
be the canonical morphism. One finds this:
\begin{center}
\makebox[0.20\linewidth][l]{(4.15.8.4)}%
\begin{minipage}[c]{0.78\linewidth}
If $u \in N_i$, then the equivalence relation $R_i$ on $M$ defined by $\alpha_i$ ($v \overset{R_i}{\sim} v' \Longleftrightarrow \alpha_i(v) = \alpha_i(v')$) is equal to that defined by the map $v \longmapsto uv$ from $M$ to $M$, i.e.
\[
\alpha_i(v) \ = \ \alpha_i(v') \qquad \Longleftrightarrow \qquad uv \ = \ uv' \ .
\]
\end{minipage}
\end{center}

This is immediate, since $I_{i,u}$ being cofinal in $I$, one has $\alpha_i(v) = \alpha_i(v')$ if and only if $\exists\ j \in I_{i,u}$ with $\alpha_{ji}v = \alpha_{ji}v'$, now $\alpha_{ji} = u$. % [sic] C'est immédiate -> C'est immédiat, agreement

For $u \in M$, let us denote by $R_u$ the equivalence relation on $M$ defined by $v \longmapsto uv$. The set of equivalence relations on \texttt{[}$M$\texttt{]}
\msp{page 455}
obtained in this way has cardinality $\leqslant \card M \leqslant \pi_0$, hence the set of equivalence relations $R_i$ ($i \in I$) also has cardinality $\leqslant \pi_0$. This shows, $\pi_0$ being large relative to $\pi_0$, that there exists an equivalence relation of the form $R_u$ ($u \in M$), and a cofinal set $I'$ in $I$, such that $\forall\ i \in I'$ one has $R_i = R_u$. This shows that, at the cost of replacing $I$ by the cofinal ordered set $I'$, that one can suppose that one has
\[
R_i \ = \ R_u \qquad \forall\, i \in I \ .
\tag*{(4.15.8.5)}
\]
Let then
\[
E \ = \ M/R_u \ = \ M/R_i \qquad (\forall\ i \in I) \ ,
\]
one therefore sees that the $\alpha_{ji} \in M$ (for $i \leqslant j$) are such that the map
\[
\tau_{\alpha_{ji}} \ : \ v \ \longmapsto \ \alpha_{ji}v \ : \ M \ \longrightarrow \ M
\]
is compatible with the equivalence relation $R_u$, hence passes to the quotient as a map
\[
\overline{\alpha_{ji}} \ : \ E \ \longmapsto \ E \ ,
\]
which is moreover \textit{injective.}

Let $M'$ be the submonoid of $M$ formed of the $u \in M$ which have the preceding stability property, as well as the injectivity property we have just stated, and $\overline{M'}$ the image monoid of $M'$ in the monoid
\[
\mathrm{Mon}(E) \ \subset \ \mathrm{End}_{\mathrm{Ens}}(E)
\]
\msp{page 456}
of injective maps from $E$ to itself. One therefore sees that the given direct system, defined by the $\alpha_{ji}$, takes its values in the submonoid $M'$ of $M$ (i.e. of the corresponding subcategory of the category $B_M$), and that the direct limit $F(\xi)$ which interests us (to prove $\card F(\xi) \leqslant \pi_0$) is isomorphic to the direct limit of the direct system of a system of sets all equal to $E$, the transition morphisms being the $\overline{\alpha_{ji}} \in \overline{M'} \subset \mathrm{Mon}(E)$. Moreover, one will have $\card \overline{M'} \leqslant \card M \leqslant \pi_0$. The desired result $\card F(\xi) \leqslant \pi_0$ will then follow from the

\textbf{Key lemma 4.15.8.1.} \textit{Let $M$ be a monoid of monomorphisms (i.e. such that $uv = uv'$ $\Longrightarrow v = v'$, for $u,v,v' \in M$), and let $(\alpha_{ji})_{i\leqslant j}$ be a direct system in $B_M$, indexed by an ordered set $I$ large relative to $\pi_0$, where $\pi_0$ is a cardinal}
\printedpage{177}
% head: A. Grothendieck : Dérivateurs, Ch. XVIII, Édition des Universités de Montpellier II et Paris VII
\msp{page 457}
\textit{which is infinite and such that} $\card M \leqslant \pi_0$. \textit{Then there exists an $i_0 \in I$ such that for every $i \geqslant i_0$, $\alpha_{i_0,i}$ is an isomorphism.}

If one applies this lemma to the case of $\overline{M'}$ and of the $\overline{\alpha_{ji}}$, one therefore finds that there exists an $i_0 \in I$ such that the $\overline{\alpha_{ji_0}}$ are isomorphisms. This implies that $F(\xi) \simeq \varinjlim\limits_{i} \underbrace{E_i}_{E}$ is isomorphic to $E_{i_0}$, hence has cardinality equal to $\card E_{i_0} = \card E \leqslant \pi_0$, which will complete the proof.

Moreover, statement 4.8.13 implies 4.15.8.1, for in the direct system of the $(M_i = M, \alpha_{ji})$, with $\varinjlim$ equal to $F(\xi)$, if there did \textit{not} exist an $i_0$ as stipulated in 4.15.8.1, then one would easily conclude from this that one would have $\card F(\xi) > \pi_0$, contrary to 4.8.13. We are going to bring this out in a lemma:

\msp{page 458}
\textbf{Lemma 4.15.8.2.} \textit{Let $(E_i)_{i\in I}$ be a direct system of sets, indexed by an ordered set $I$ large relative to $\pi_0$. We suppose that there exists $i_0 \in I$ such that}
\begin{enumerate}[nosep]
\item[\textit{a)}] \textit{$i \geqslant i_0$ implies} $\card E_i = \pi_0$, \textit{and}
\item[\textit{b)}] \textit{$j \geqslant i \geqslant i_0$ implies that $u_{ji} : E_i \longrightarrow E_j$ is a monomorphism. Suppose moreover}
\item[\textit{c)}] \textit{$\forall\ i \geqslant i_0$ $\exists\ j > i$ such that $u_{ji}$ is} not \textit{an isomorphism (i.e. is not surjective). Let $E_\infty = \varinjlim\limits_{i\in I} E_i$. Then} $\card E_\infty > \pi_0$.
\end{enumerate}
This lemma therefore implies the \textit{equivalence} of 4.8.13 and the ``key lemma'' 4.15.8.1. To prove 4.15.8.2, we are going to prove something a little more precise:

\textbf{Corollary 4.15.8.3.} \textit{Under the conditions of 4.15.8.2, let $\alpha_0$ be the first ordinal such that} $\card \alpha_0 > \pi_0$, \textit{and let $I_{\alpha_0}$ be the set of ordinals $\beta < \alpha$ (an ordered set which is large relative to $\pi_0$). There then exists a strictly increasing map}
\[
I_{\alpha_0} \ \longrightarrow \ I \ , \qquad \alpha \ \longmapsto \ i_\alpha \ , \qquad \textit{mapping } 0 \textit{ to } i_0 \ ,
\]
\msp{page 459}
\textit{such that for every $\alpha \in I_{\alpha_0}$, one has}
\[
\underbrace{u_{i_{\alpha+1},i_\alpha}}_{\substack{\textit{it is a } \text{monomorphism,}\\ \textit{by b, since } i_{\alpha+1} > i_\alpha \geqslant i_0}} \qquad\qquad \text{\textit{is not an isomorphism} (\footnotemark[187]).}
\]
\footnotetext[187]{\textit{better}: $\forall\ \alpha \in I_{\alpha_0}$ $\varinjlim\limits_{\beta\in I_\alpha} E_{i_\beta} \longrightarrow E_{i_\alpha}$ is not surjective.}
\textit{One then has}
\[
E'_\infty \ \overset{\text{def}}{=} \ \underbrace{\varinjlim_{\alpha\in I_{\alpha_0}} E_{i_\alpha}}_{} \ \hookrightarrow \ \varinjlim_{i\in I} E_i \ \overset{\text{def}}{=} E_\infty
\]
\textit{and}
\[
\card E'_\alpha \ \geqslant \ \underbrace{\card I_{\alpha_0}}_{=\ \textit{successor of } \pi_0} \ > \ \pi_0 \ ,
\] % [sic] E'_\alpha -> E'_\infty, the following line has \card E'_\infty
\textit{whence} $\card E_\infty \geqslant \card E'_\infty > \pi_0$.

\textsc{Proof}. The construction of the $i_\alpha$ by transfinite induction is immediate, using c) and the fact that $I$ is large relative to $\pi_0$, and that for every $\alpha \in I_{\alpha_0}$, the set $I_\alpha$
\printedpage{178}
% head: A. Grothendieck : Dérivateurs, Ch. XVIII, Édition des Universités de Montpellier II et Paris VII
of $\beta < \alpha$ has cardinality $\leqslant \pi_0$. The fact that $E'_\infty \longrightarrow E_\infty$ is injective follows from b). The fact that
\[
\card E'_\infty \ \geqslant \ \card I_{\alpha_0}
\]
follows from the fact that $E'_\infty$ is an increasing union of subsets identified with $E_{i_\alpha}$, and for every
\msp{page 460}
$\alpha$ one can find a
\[
\xi_\alpha \ \in \ E_{i_\alpha}\backslash \bigcup_{\beta<\alpha} E_{i_\beta} \ .
\]
The map
\[
\begin{aligned}
\alpha \ &\longmapsto \ \xi_\alpha\\
I_{\alpha_0} \ &\longrightarrow \ E'_\infty
\end{aligned}
\]
is injective, hence $\card E'_\infty \geqslant \card I_{\alpha_0}$.

% vertical bar printed in the left margin beside the whole of Corollaire 4.15.8.4
\textbf{Corollary 4.15.8.4.} \textit{For given $M$ and $\pi_0$, for the key lemma 4.15.8.1 to be valid, it is necessary and sufficient that it be so in the case where $I = I_{\alpha_0}$, $\alpha_0$ being the first ordinal such that} $\card \alpha_0 > \alpha_0$. \textit{It amounts to the same thing that one cannot find a direct system $(\alpha_{ji})_{j\geqslant i}$ indexed by $I_{\alpha_0}$, with values in $M$, such that for $i < j$, one has $\alpha_{ji}$ noninvertible in $M$.} % [sic] \card \alpha_0 > \alpha_0 -> \card \alpha_0 > \pi_0, cf. Corollaire 4.15.8.3

\msp{page 461}
I am going to summarize what has been obtained in a summary statement:

\textbf{Theorem 4.15.9.} \textit{Let $\mathcal{C}$ be a small category, $\pi$ a cardinal such that $\pi \geqslant$} $\card\mathcal{C}$. \textit{The following conditions are equivalent (and I conjecture that they are always satisfied):}
\begin{enumerate}[nosep]
\item[\textit{a)}] \textit{$\mathcal{M} \overset{\text{def}}{=} \mathrm{Ind}_\pi\mathcal{C}$ is reduced to $\mathcal{M}_\pi$.}
\item[\textit{a$'$)}] \textit{(If $\mathcal{C}$ is Karoubian.) Every direct system $\underline{x} = (x_i)_{i\in I}$ in $\mathcal{C}$, large relative to $\pi$, is essentially constant.}
\item[\textit{a$''$)}] \textit{For every direct system $(x_i)_{i\in I}$ in $\mathcal{C}$, large relative to $\pi$, there exists a cofinal functor $J \longrightarrow I$, with $J$ an ordered set large relative to $\pi$, and such that the reindexed system $(x_j)_{j\in J}$ satisfies the following conditions:}
\begin{enumerate}[nosep]
\item[\textit{1)}] \textit{All the $x_j$ are equal to one and the same element $x$ of $\mathcal{C}$.}
\item[\textit{2)}] \textit{There exists a projector $p$ in $x$, such that one has $\alpha_{j'j} = p$ for $j < j'$ ($j,j' \in J$), where $\alpha_{j'j} : x_j \longrightarrow x_{j'}$ is the transition arrow.}
\end{enumerate}
\item[\textit{b)}] \textit{For every $x \in \mathrm{Ob}\,\mathcal{C}$, and for every quotient monoid $M'$ of $M = \mathrm{End}_{\mathcal{C}}(x)$, such that $M'$ is a monoid of monomorphisms (i.e. $uv = uv' \Longrightarrow v = v'$ if $u,v,v' \in M'$), and for every direct system large relative to $\pi$ with values in $B_{M'}$ (defined by a system of elements $(\alpha_{ji})_{i\leqslant j,\ i,j\in I}$ of an ordered set $I$ large relative to $\pi$, with $\alpha_{kj}\alpha_{ji} = \alpha_{ki}$ for $i \leqslant j \leqslant k$), there exists $i_0 \in I$ such}
\msp{page 462}
\textit{that $\alpha_{ii_0} = 1$ for $i \geqslant i_0$, i.e. $\alpha_{ji} = 1$ for $j \geqslant i \geqslant i_0$.}
\item[\textit{b$'$)}] \textit{$M'$ being a monoid of monomorphisms as in b) (hence} $\card M' \leqslant \card\mathcal{C} \leqslant$ \textit{$\pi$), let moreover $\alpha_0$ be the smallest ordinal such that} $\card\alpha_0 > \pi$, \textit{and let $I = I_{\alpha_0}$ be the ordered set of ordinals $\alpha < \alpha_0$ (which set is large relative to $\pi$). There exists} no \textit{direct system in $B_{M'}$, such that for $i < j$ in $I$, $\alpha_{ji}$ is a noninvertible element in $M'$ (hence $u \longmapsto \alpha_{ji}u : M' \longrightarrow M'$ is injective but not surjective).}
\end{enumerate}

\end{document}
